\section{Numerical Methodology}

\subsection{Discretization and Grid Strategy}
The numerical evaluation of the emission integrals requires a careful balance between resolution and computational cost. We employ a discretized energy grid spanning $[0, 12]$~eV with approximately $10^3$ points. This density is sufficient to resolve the sharp Fermi edge at low temperatures and the relevant photon energy range for visible and near-infrared emission. For momentum-space calculations, we utilize a cylindrical coordinate system $(k_{\perp}, k_{\parallel})$ aligned with the relevant valley axes. The grid dimensions are tuned dynamically based on the convergence behavior of the delta function approximation.

\subsection{Gaussian Delta Approximation}
\label{sec:delta}
Central to our numerical approach is the regularized representation of the Dirac delta function $\delta(x)$ as a Gaussian distribution $G_{\sigma}(x)$. This introduces two control parameters: the width $\sigma$, which determines the approximation error, and the sampling density, which determines the quadrature error.

We adhere to a strict convergence protocol derived from our benchmarks. For a given $\sigma$, we define an integration window of width $\pm m\sigma$ (typically $m=10$) to capture the Gaussian tail. Within this window, the grid step $\Delta x$ is set to a fraction $s\sigma$ (typically $s=0.1$) to ensure at least $2m/s \approx 200$ sample points across the peak. Convergence is declared when the relative change in the integral value upon further reducing $\sigma$ falls below a target tolerance of $10^{-3}$. In anisotropic scenarios, we perform the coordinate transformation to isotropic space *before* applying the cutoff, ensuring that the Gaussian constraint acts uniformly on the radial coordinate $q^2$.

\subsection{Implementation Stack}
The core integration logic is implemented in a series of Jupyter notebooks utilizing the Python/NumPy/Matplotlib stack. To ensure reproducibility, the project runs in a dedicated local virtual environment, isolating the dependencies. The key notebooks are partitioned by physics module: \texttt{intra\_energy} for the 1D energy-space benchmarks, \texttt{intra\_momentum} for the high-dimensional convergence tests, and \texttt{inter\_momentum} for the anisotropic valley calculations.
